% Slide 14 — Geração de rótulos e PSE
\begin{frame}{Geração de rótulos e o PSE}
  \[
    \hat{\pi}^*(i) = \arg\min_{\pi \in \mathcal{P}} \; \mathbb{E}[\mathrm{CTI}(\pi,i)]
    \quad \text{sujeito a } \mathbb{E}[\mathrm{NS}(\pi,i)] \geq 0{,}70
  \]
  \small Para cada série, o rótulo é a política viável de \highlight{menor
  CTI} entre as que atingem o NS mínimo.
  \vspace{0.45em}

  \begin{columns}[T,onlytextwidth]
    \column{0.56\textwidth}
      \textbf{PSE (estágio atual)}
      \small
      \begin{itemize}
        \item Classificador XGBoost multiclasse
        \item Entrada: vetor de 17 características
        \item Saída: política recomendada + confiança + \textit{ranking}
      \end{itemize}
    \column{0.44\textwidth}
      \begin{alertblock}{Ainda não finalizado}
        \footnotesize O PSE foi formulado e testado nos rótulos disponíveis; o
        treinamento em escala e a validação de generalização são etapa da
        expansão nacional.
      \end{alertblock}
  \end{columns}
  \note{
    O rótulo de cada série vem diretamente do benchmark simulado: entre as
    políticas que atingem o nível de serviço mínimo de 0,70, escolhe-se a
    de menor custo total de inventário. Esses pares -- características e
    rótulo -- formam o dataset supervisionado do PSE, que no estágio atual
    é um classificador XGBoost. É importante deixar claro que o PSE não
    está finalizado: ele foi formulado e testado sobre os rótulos
    disponíveis dos dois experimentos, mas o treinamento em escala nacional
    e a validação de generalização para novas séries e regimes ainda são
    etapas futuras.
  }
\end{frame}
